% =====================================================================
\chapter{Teoremas del valor medio y aplicaciones}
\label{tema:5}
% =====================================================================

La derivada es una información local: describe $f$ en un entorno infinitesimal de cada punto. Los teoremas de este tema convierten esa información local en información global sobre un intervalo: si la derivada es positiva en todo un intervalo, la función crece en él; si está acotada, la función no puede variar más que una cantidad proporcional a la longitud del intervalo. El instrumento es el teorema del valor medio, del que se derivan la regla de L'Hôpital, los criterios de extremos y la teoría de la convexidad. Se siguen el capítulo~11 de Spivak~\cite{spivak}, el capítulo~4 de Apostol~\cite{apostol}, la §5.3 de los apuntes de \emph{Analysis~I}~\cite{cambA1}, las §4.1--4.5 y~4.7 de Stewart~\cite{stewart} y la unidad~2 del curso 18.01SC del MIT~\cite{mit1801}.

\section{Extremos locales y el teorema de Rolle}

\begin{definicion}[Extremos]\label{def:extremos}
Sea $f\colon D\to\R$ y $c\in D$. Se dice que $f$ tiene en $c$ un \emph{máximo absoluto} (o global) si $f(x)\le f(c)$ para todo $x\in D$, y un \emph{máximo local} (o relativo) si existe $\delta>0$ con $f(x)\le f(c)$ para todo $x\in D$ con $\abs{x-c}<\delta$. Análogamente se definen los mínimos. Un punto $c$ interior a $D$ con $f'(c)=0$ se llama \emph{punto crítico} o estacionario.
\end{definicion}

\begin{teorema}[Fermat: condición necesaria de extremo]\label{teo:fermat}
Si $f$ tiene un extremo local en un punto $c$ interior a su dominio y es derivable en $c$, entonces $f'(c)=0$.
\end{teorema}

El recíproco es falso: $x^{3}$ tiene $f'(0)=0$ y no tiene extremo en~0. Tampoco se puede omitir la hipótesis de derivabilidad ($\abs x$ tiene un mínimo en~0) ni la de punto interior ($x$ en $[0,1]$ tiene su máximo en~1, con $f'(1)=1$).

\begin{teorema}[Rolle]\label{teo:rolle}
Sea $f$ continua en $[a,b]$ y derivable en $(a,b)$, con $f(a)=f(b)$. Entonces existe $c\in(a,b)$ con $f'(c)=0$.
\end{teorema}

\section{El teorema del valor medio}

\begin{teorema}[Teorema del valor medio de Lagrange]\label{teo:tvm}
Sea $f$ continua en $[a,b]$ y derivable en $(a,b)$. Entonces existe $c\in(a,b)$ con
\[
f'(c)=\frac{f(b)-f(a)}{b-a},\qquad\text{es decir,}\qquad f(b)-f(a)=f'(c)(b-a).
\]
\end{teorema}

Geométricamente: hay un punto en el que la tangente es paralela a la cuerda. Cinemáticamente: si un coche recorre 121 kilómetros en una hora y cincuenta minutos, su velocidad media es de 66~km/h, y en algún instante el velocímetro marca exactamente esa velocidad.

\begin{corolario}[Consecuencias del teorema del valor medio]\label{cor:tvm}
Sea $f$ continua en un intervalo $I$ y derivable en su interior.
\begin{enumerate}
  \item Si $f'=0$ en el interior de $I$, $f$ es constante en $I$.
  \item Si $f'=g'$ en el interior de $I$, entonces $f-g$ es constante.
  \item Si $f'>0$ (respectivamente $f'\ge0$) en el interior de $I$, $f$ es estrictamente creciente (respectivamente creciente) en $I$; análogamente con $<0$ y $\le0$.
  \item Si $\abs{f'(x)}\le M$ en el interior de $I$, entonces $\abs{f(x)-f(y)}\le M\abs{x-y}$ para todos $x,y\in I$; se dice que $f$ es \emph{lipschitziana} de constante~$M$.
\end{enumerate}
\end{corolario}

\addtocounter{definicion}{1}%

\begin{teorema}[Teorema del valor medio de Cauchy]\label{teo:cauchy}
Sean $f$ y $g$ continuas en $[a,b]$ y derivables en $(a,b)$, con $g'(x)\ne0$ para todo $x\in(a,b)$. Entonces $g(b)\ne g(a)$ y existe $c\in(a,b)$ con
\[
\frac{f'(c)}{g'(c)}=\frac{f(b)-f(a)}{g(b)-g(a)}.
\]
\end{teorema}

\section{La regla de L'Hôpital}

\begin{teorema}[Regla de L'Hôpital]\label{teo:lhopital}
Sean $f$ y $g$ derivables en un intervalo abierto que contiene a $a$, salvo quizá en $a$, con $g'(x)\ne0$ para $x\ne a$, y supóngase que $\lim_{x\to a}f(x)=\lim_{x\to a}g(x)=0$. Si existe $\lim_{x\to a}\dfrac{f'(x)}{g'(x)}=\ell$ (finito o infinito), entonces
\[
\lim_{x\to a}\frac{f(x)}{g(x)}=\ell.
\]
El resultado vale también para límites laterales, para $a=\pm\infty$, y cuando $\lim f=\lim g=\pm\infty$ (forma $\infty/\infty$).
\end{teorema}

\addtocounter{definicion}{1}%

\begin{proposicion}[Jerarquía de crecimientos]\label{prop:crecimientos}
Para todos $p>0$ y $b>1$: $\displaystyle\lim_{x\to+\infty}\frac{\ln x}{x^{p}}=0$ y $\displaystyle\lim_{x\to+\infty}\frac{x^{p}}{b^{x}}=0$.
\end{proposicion}

\section{Monotonía, extremos y optimización}

\begin{proposicion}[Criterios de extremo local]\label{prop:criterios}
Sea $f$ derivable en un entorno de un punto crítico $c$.
\begin{enumerate}
  \item (\emph{Criterio de la primera derivada}) Si $f'\ge0$ a la izquierda de $c$ y $f'\le0$ a la derecha (en un entorno), $f$ tiene un máximo local en $c$; con los signos opuestos, un mínimo local.
  \item (\emph{Criterio de la segunda derivada}) Si $f''(c)$ existe y $f''(c)<0$, $f$ tiene un máximo local estricto en $c$; si $f''(c)>0$, un mínimo local estricto. Si $f''(c)=0$, el criterio no decide.
\end{enumerate}
\end{proposicion}

\begin{proposicion}[Extremos absolutos en un intervalo cerrado]\label{prop:intervalo-cerrado}
Sea $f$ continua en $[a,b]$. Su máximo y su mínimo absolutos (que existen por el \cref{teo:weierstrass}) se alcanzan en los extremos $a$, $b$, en puntos críticos de $(a,b)$ o en puntos de $(a,b)$ donde $f$ no es derivable. Basta, pues, evaluar $f$ en esos candidatos y comparar.
\end{proposicion}

\addtocounter{definicion}{1}%

\section{Convexidad}

\begin{definicion}[Función convexa]\label{def:convexa}
Una función $f\colon I\to\R$ en un intervalo $I$ es \emph{convexa} si para todos $x,y\in I$ y todo $t\in[0,1]$
\[
f\bigl(tx+(1-t)y\bigr)\le tf(x)+(1-t)f(y),
\]
es decir, si el segmento que une dos puntos de la gráfica queda por encima de ella. Es \emph{estrictamente convexa} si la desigualdad es estricta para $x\ne y$ y $t\in(0,1)$, y \emph{cóncava} si $-f$ es convexa. Un \emph{punto de inflexión} es un punto donde $f$ es continua y pasa de convexa a cóncava o al revés.
\end{definicion}

\begin{teorema}[Caracterizaciones de la convexidad]\label{teo:convexidad}
Sea $f$ derivable en un intervalo abierto $I$.
\begin{enumerate}
  \item Si $f'$ es creciente en $I$, entonces $f$ está por encima de sus tangentes: $f(y)\ge f(x)+f'(x)(y-x)$ para todos $x,y\in I$.
  \item Si $f$ está por encima de sus tangentes, $f$ es convexa.
  \item En particular, si $f''\ge0$ en $I$, $f$ es convexa; si $f''>0$, es estrictamente convexa.
\end{enumerate}
\end{teorema}

\addtocounter{definicion}{1}%

\addtocounter{definicion}{1}%

\addtocounter{definicion}{1}%

